\section{Conclusion}

This route solves the problem through three connected layers. First, a multi-indicator evaluation model converts heterogeneous supplier information into a stable candidate structure. Second, the candidate set is transformed into an executable planning model that produces period-by-period decisions under explicit demand and capacity constraints. Third, feasibility audits and scenario comparisons are used to determine whether the recommendation remains credible under disturbance.

The most important conclusion is that the paper no longer stops at ranking. Instead, the ranking result is treated as a planning input, and the final recommendation is expressed through an executable decision table together with a feasibility judgment. This is precisely the difference between an appearance-level draft and a route-complete paper.

In summary, the proposed route not only identifies promising suppliers, but also yields an interpretable and auditable operational plan. Within the tested scenario range, the recommendation can be treated as stable enough for research review, while the remaining uncertainty should still be handled by later human verification and stronger real-data experiments.
